% Wersja źródłowa: audyt II, jednostka 13
\setcounter{section}{12}
\section{Przestrzenie Hilberta: geometria, reprezentacja Riesza i formy półtoraliniowe}
\label{sec:hilbert13}

\subsection{Iloczyn skalarny i zupełność}

\begin{definicja}[Iloczyn skalarny i przestrzeń Hilberta]
\label{def:iloczyn-skalarny13}
Niech $H$ będzie przestrzenią wektorową nad $\K\in\{\R,\C\}$.
\emph{Iloczynem skalarnym} na $H$ nazywamy odwzorowanie
$\langle\cdot,\cdot\rangle:H\times H\to\K$ spełniające, dla
$x,y,z\in H$ i $\alpha,\beta\in\K$:
$\langle\alpha x+\beta y,z\rangle=\alpha\langle x,z\rangle+
\beta\langle y,z\rangle$,
$\langle x,y\rangle=\overline{\langle y,x\rangle}$ oraz
$\langle x,x\rangle>0$ dla $x\ne0$.
\emph{W całym rozdziale iloczyn skalarny jest liniowy w pierwszym
argumencie} i antyliniowy w drugim. Przestrzenią Hilberta nazywamy
przestrzeń z iloczynem skalarnym zupełną w normie
$\|x\|=\sqrt{\langle x,x\rangle}$.
\end{definicja}

\begin{lemat}[Nierówność Cauchy'ego--Schwarza]
\label{lem:cs13}
Dla dowolnych $x,y$ w przestrzeni z iloczynem skalarnym zachodzi
$|\langle x,y\rangle|\le\|x\|\|y\|$.
Równość dla $y\ne0$ zachodzi dokładnie wtedy, gdy $x$ jest
wielokrotnością $y$.
\end{lemat}
\begin{proof}
Dla $y\ne0$ przyjmijmy $\alpha=\langle x,y\rangle/\|y\|^2$.
Wówczas $\langle x-\alpha y,y\rangle=0$ oraz
$0\le\|x-\alpha y\|^2=\|x\|^2-|\langle x,y\rangle|^2/\|y\|^2$.
Równość oznacza $x=\alpha y$. Dla $y=0$ teza jest oczywista.
Nierówność trójkąta dla normy wynika teraz z rozwinięcia
$\|x+y\|^2$ i nierówności Cauchy'ego--Schwarza.
\end{proof}

\begin{stwierdzenie}[Tożsamość równoległoboku i polaryzacja]
\label{prop:parallelogram13}
W każdej przestrzeni z iloczynem skalarnym zachodzi
$\|x+y\|^2+\|x-y\|^2=2\|x\|^2+2\|y\|^2$.
Dla $\K=\C$ iloczyn skalarny jest odtworzony z normy wzorem
\[
\langle x,y\rangle=\frac14\bigl(\|x+y\|^2-\|x-y\|^2
+i\|x+iy\|^2-i\|x-iy\|^2\bigr).
\]
Dla $\K=\R$ wystarczą pierwsze dwa wyrazy.
\end{stwierdzenie}
\begin{proof}
Rozwijamy kwadraty norm. Przy naszej konwencji
$\|x+iy\|^2-\|x-iy\|^2=4\operatorname{Im}\langle x,y\rangle$;
różnica dla $x\pm y$ to $4\operatorname{Re}\langle x,y\rangle$.
\end{proof}

\begin{przyklad}[Przestrzenie $\ell^2$ i $L^2$]
\label{ex:hilbert-l2-13}
Na $\ell^2$ przyjmujemy $\langle x,y\rangle=\sum_{n\ge1}x_n\overline{y_n}$.
Na $L^2(\Omega,\mu)$ przyjmujemy
$\langle f,g\rangle=\int_\Omega f(x)\overline{g(x)}\,\mathrm d\mu(x)$.
Nierówność Cauchy'ego--Schwarza zapewnia zbieżność i ciągłość tych
wyrażeń, zaś zupełność wynika z
twierdzeń o zupełności przestrzeni $\ell^2$ i $L^2$
(drugi przykład przy założeniach o mierze z \hyperref[sec:klasyczne-banach12]{rozdziału~12}).
Przestrzeń $C([0,1])$ z normą supremum nie jest przestrzenią Hilberta w tej normie:
dla $f(t)=1$ i $g(t)=t$ mamy
$\|f+g\|_\infty^2+\|f-g\|_\infty^2=5$,
podczas gdy $2\|f\|_\infty^2+2\|g\|_\infty^2=4$.
\end{przyklad}

\subsection{Rzut na zbiór wypukły i ortogonalność}

\begin{twierdzenie}[Twierdzenie o rzucie metrycznym]
\label{thm:projection-convex13}
Niech $H$ będzie rzeczywistą lub zespoloną przestrzenią Hilberta,
a $C\subseteq H$ niepustym, domkniętym i wypukłym zbiorem.
Dla każdego $x\in H$ istnieje dokładnie jeden punkt $p=P_Cx\in C$
spełniający $\|x-p\|=\inf_{z\in C}\|x-z\|$.
Punkt $p\in C$ jest tym minimizatorem wtedy i tylko wtedy, gdy
\[
\operatorname{Re}\langle x-p,z-p\rangle\le0\qquad(z\in C).
\]
Ponadto odwzorowanie $P_C:H\to C$ jest $1$-lipschitzowskie.
\end{twierdzenie}
\begin{proof}
Niech $d=\inf_{z\in C}\|x-z\|$ i wybierzmy $z_n\in C$ z
$\|x-z_n\|^2\to d^2$. Ponieważ $(z_n+z_m)/2\in C$,
tożsamość równoległoboku zastosowana do $x-z_n$ i $x-z_m$ daje
\[
\|z_n-z_m\|^2
=2\|x-z_n\|^2+2\|x-z_m\|^2
-4\left\|x-\frac{z_n+z_m}{2}\right\|^2
\le2\|x-z_n\|^2+2\|x-z_m\|^2-4d^2\longrightarrow0.
\]
Zatem $(z_n)$ jest ciągiem Cauchy'ego. Zupełność $H$ i domkniętość
$C$ dają granicę $p\in C$, a ciągłość normy daje $\|x-p\|=d$.
Jeśli $p,q$ są minimizatorami, ten sam wzór daje $p=q$.

Niech $z\in C$ i $t\in[0,1]$. Minimalność $p$ oznacza
$\|x-p\|^2\le\|x-(p+t(z-p))\|^2$,
a więc
$0\le-2t\operatorname{Re}\langle x-p,z-p\rangle+t^2\|z-p\|^2$.
Dzieląc przez $t>0$ i przechodząc z $t$ do zera otrzymujemy
warunek wariacyjny. Odwrotnie, warunek ten i rozwinięcie
$\|x-z\|^2=\|x-p\|^2+\|z-p\|^2-
2\operatorname{Re}\langle x-p,z-p\rangle$ implikują minimalność.

Dla $p=P_Cx$, $q=P_Cy$ zastosujmy warunek kolejno z $z=q$
i $z=p$. Otrzymujemy
$\operatorname{Re}\langle x-p,q-p\rangle\le0$ oraz
$\operatorname{Re}\langle y-q,p-q\rangle\le0$. Po dodaniu
$\|p-q\|^2\le\operatorname{Re}\langle x-y,p-q\rangle
\le\|x-y\|\|p-q\|$, skąd $\|p-q\|\le\|x-y\|$.
\end{proof}

\begin{definicja}[Dopełnienie ortogonalne]
\label{def:orth-complement13}
Dla $M\subseteq H$ definiujemy
$M^\perp=\{x\in H:\langle x,m\rangle=0\text{ dla wszystkich }m\in M\}$.
Jest to domknięta podprzestrzeń liniowa $H$; istotnie, to
przecięcie jąder ciągłych funkcjonałów $x\mapsto\langle x,m\rangle$.
\end{definicja}

\begin{twierdzenie}[Rozkład ortogonalny]
\label{thm:orth-decomposition13}
Niech $H$ będzie przestrzenią Hilberta, a $M\subseteq H$ domkniętą
podprzestrzenią liniową. Wówczas dla każdego $x\in H$ istnieje
jedyny rozkład $x=m+n$, $m\in M$, $n\in M^\perp$.
Odwzorowanie $P_Mx=m$ jest liniowym operatorem o normie co najwyżej
$1$, spełnia $P_M^2=P_M=P_M^*$ oraz
$\|x\|^2=\|P_Mx\|^2+\|x-P_Mx\|^2$.
Ponadto $M^{\perp\perp}=M$.
Jeśli $M\ne\{0\}$, to $\|P_M\|=1$.
\end{twierdzenie}
\begin{proof}
Z Twierdzenia~\ref{thm:projection-convex13} istnieje $m=P_Mx$.
Warunek wariacyjny zastosowany do $m+tv$ dla każdego $v\in M$
i $t\in\R$, a dla przestrzeni zespolonej także do $m+itv$,
daje $\langle x-m,v\rangle=0$. Stąd $n=x-m\in M^\perp$.
Jeśli $x=m_1+n_1=m_2+n_2$, to $m_1-m_2=n_2-n_1$ należy
do $M\cap M^\perp=\{0\}$. Jednoznaczność rozkładu zapewnia
liniowość $P_M$; twierdzenie Pitagorasa daje jego ograniczoność.

Dla $x=m+n$ i $y=u+v$ mamy
$\langle P_Mx,y\rangle=\langle m,u\rangle
=\langle x,P_My\rangle$; to oznacza $P_M=P_M^*$,
przy czym istnienie operatora sprzężonego wykażemy formalnie
w Twierdzeniu~\ref{thm:adjoint13} poniżej.
Z rozkładu $H=M\oplus M^\perp$ wynika, że
$M^{\perp\perp}=M$. Dla $m\in M\setminus\{0\}$ zachodzi
$P_Mm=m$, zatem norma operatorowa nie może być mniejsza niż $1$.
\end{proof}

\begin{przyklad}[Rzut na funkcje stałe]
\label{ex:projection-mean13}
W przestrzeni $H=L^2([0,1],\lambda)$ niech
$M=\{\text{klasy funkcji stałych}\}$. Wtedy
$P_Mf=(\int_{[0,1]}f(t)\,\mathrm d\lambda(t))\chi_{[0,1]}$.
Rzeczywiście, $\langle f-P_Mf,\chi_{[0,1]}\rangle=0$.
Ogólniej, rzut na linię $\K u$, $u\ne0$, ma wzór
$P_{\K u}x=\langle x,u\rangle u/\|u\|^2$.
\end{przyklad}

\subsection{Twierdzenie Riesza o reprezentacji funkcjonałów}

\begin{twierdzenie}[Twierdzenie Riesza dla przestrzeni Hilberta]
\label{thm:riesz-hilbert13}
Niech $H$ będzie przestrzenią Hilberta nad $\K\in\{\R,\C\}$.
Dla każdego ciągłego $\K$-liniowego funkcjonału $F:H\to\K$
istnieje dokładnie jeden $y\in H$ taki, że
$F(x)=\langle x,y\rangle$ dla wszystkich $x\in H$.
Ponadto $\|F\|_{H'}=\|y\|$.
Odwzorowanie $J:H\to H'$, $J(y)(x)=\langle x,y\rangle$,
jest izometrycznym izomorfizmem \emph{antyliniowym} dla $\K=\C$
i liniowym dla $\K=\R$.
\end{twierdzenie}
\begin{proof}
Jeżeli $F=0$, wybieramy $y=0$. Załóżmy $F\ne0$.
Podprzestrzeń $M=\ker F$ jest domknięta, a $M^\perp\ne\{0\}$:
gdyby $M^\perp=\{0\}$, to z Twierdzenia~\ref{thm:orth-decomposition13}
wynikałoby $M=H$.
Niech $z\in M^\perp$, $z\ne0$. Wtedy $F(z)\ne0$, gdyż
$z\in M\cap M^\perp$ implikowałoby $z=0$.
Dla dowolnego $x\in H$ mamy
$x-F(x)z/F(z)\in M$, więc
\[
0=\left\langle x-\frac{F(x)}{F(z)}z,z\right\rangle
=\langle x,z\rangle-\frac{F(x)}{F(z)}\|z\|^2.
\]
Zatem $F(x)=F(z)\langle x,z\rangle/\|z\|^2$.
Przy naszej konwencji należy wziąć
$y=\overline{F(z)}z/\|z\|^2$, ponieważ
$\langle x,cy\rangle=\overline c\langle x,y\rangle$.
Nierówność Cauchy'ego--Schwarza daje $\|F\|\le\|y\|$,
a podstawienie $x=y/\|y\|$ dla $y\ne0$ daje równość.
Jeśli $\langle x,y_1-y_2\rangle=0$ dla wszystkich $x$, to
$x=y_1-y_2$ zapewnia $y_1=y_2$. Antyliniowość $J$
wynika bezpośrednio z $J(\alpha y+\beta z)
=\overline\alpha J(y)+\overline\beta J(z)$, ponieważ
iloczyn skalarny jest antyliniowy w drugim argumencie.
\end{proof}

\begin{uwaga}[Dwa różne twierdzenia Riesza]
\label{rem:two-riesz13}
Twierdzenie~\ref{thm:riesz-hilbert13} identyfikuje przestrzeń sprzężoną
przestrzeni Hilberta z nią samą (antyliniowo w przypadku zespolonym).
Natomiast twierdzenie Riesza--Markowa z teorii przestrzeni $C(K)$ identyfikuje
przestrzeń sprzężoną $C(K)$ z przestrzenią miar. Nie są to
odmienne wersje tej samej reprezentacji w sensie konstrukcyjnym.
\end{uwaga}

\subsection{Rodziny ortonormalne i bazy przestrzeni Hilberta}

\begin{definicja}[Rodzina ortonormalna i baza przestrzeni Hilberta]
\label{def:on-basis13}
Niech $H$ będzie przestrzenią Hilberta, a $I$ zbiorem indeksów.
Rodzina $(e_i)_{i\in I}\subseteq H$ jest \emph{ortonormalna}, jeśli
$\langle e_i,e_j\rangle=\delta_{ij}$.
Nazywamy ją \emph{bazą ortonormalną przestrzeni Hilberta}, jeśli
$\overline{\lin\{e_i:i\in I\}}=H$.
Nie należy utożsamiać bazy ortonormalnej przestrzeni Hilberta z algebraiczną bazą Hamela.
\end{definicja}

\begin{lemat}[Bessel i ortogonalne sumy skończone]
\label{lem:bessel13}
Niech $H$ będzie przestrzenią Hilberta, $(e_i)_{i\in I}$ rodziną
ortonormalną, a $x\in H$. Dla każdego skończonego $A\subseteq I$
zachodzi
\[
\left\|x-\sum_{i\in A}\langle x,e_i\rangle e_i\right\|^2
=\|x\|^2-\sum_{i\in A}|\langle x,e_i\rangle|^2.
\]
W szczególności $\sum_{i\in A}|\langle x,e_i\rangle|^2\le\|x\|^2$
dla każdego takiego $A$ (nierówność Bessela).
\end{lemat}
\begin{proof}
Wektor $x-\sum_{i\in A}\langle x,e_i\rangle e_i$ jest
prostopadły do każdego $e_j$, $j\in A$. Z twierdzenia Pitagorasa
dla jego rozkładu z sumą ortogonalną wynika wzór.
\end{proof}


\begin{twierdzenie}[Przeliczalność nośnika współczynników Fouriera]
\label{thm:countable-fourier-support13}
Niech $H$ będzie przestrzenią Hilberta nad $\K\in\{\R,\C\}$,
$(e_i)_{i\in I}$ dowolną rodziną ortonormalną (zbiór indeksów $I$
może być nieprzeliczalny), a $x\in H$. Wówczas zbiór
$ I_x:=\{i\in I:\langle x,e_i\rangle\ne0\}$
jest co najwyżej przeliczalny. Dokładniej, dla każdego $n\in\N$,
$n\ge1$, zbiór
$ I_{x,n}:=\{i\in I:|\langle x,e_i\rangle|\ge1/n\}$
jest skończony i ma co najwyżej $n^2\|x\|^2$ elementów.
\end{twierdzenie}
\begin{proof}
Niech $F\subseteq I_{x,n}$ będzie zbiorem skończonym.
Z nierówności Bessela (Lemat~\ref{lem:bessel13}) otrzymujemy
$|F|/n^2\le\sum_{i\in F}|\langle x,e_i\rangle|^2\le\|x\|^2$.
Gdyby $I_{x,n}$ miał więcej niż $n^2\|x\|^2$ elementów,
można byłoby wybrać skończony podzbiór o liczności
przekraczającej tę liczbę, co przeczyłoby otrzymanemu oszacowaniu.
Każdy $I_{x,n}$ jest zatem skończony. Dla każdego
$i\in I_x$ liczba $|\langle x,e_i\rangle|$ jest dodatnia, więc
przekracza $1/n$ dla pewnego $n\ge1$. Stąd
$I_x=\bigcup_{n\ge1}I_{x,n}$ jest sumą przeliczalnej rodziny
zbiorów skończonych.
\end{proof}

\begin{uwaga}[Dlaczego warunek jest silniejszy od samego Bessela?]
\label{rem:fourier-countability13}
Nie zakładaliśmy ośrodkowości $H$ ani przeliczalności $I$.
Twierdzenie dotyczy \emph{każdego pojedynczego wektora}:
przy nieprzeliczalnej bazie ortonormalnej różne wektory mogą
mieć różne przeliczalne nośniki współczynników. Nie wynika
stąd, że całą bazę można ponumerować jednym ciągiem.
\end{uwaga}

\begin{twierdzenie}[Rozwinięcie Fouriera, tożsamość Parsevala i zbieżność bezwarunkowa]
\label{thm:parseval13}
\label{cor:fourier-enumeration13}
Niech $H$ będzie przestrzenią Hilberta nad $\K\in\{\R,\C\}$,
$(e_i)_{i\in I}$ rodziną ortonormalną, a
$M=\overline{\lin\{e_i:i\in I\}}$. Dla każdego $x\in H$
sieć sum $s_A(x)=\sum_{i\in A}\langle x,e_i\rangle e_i$,
indeksowana skończonymi podzbiorami $A\subseteq I$ uporządkowanymi
przez inkluzję, zbiega w normie do $P_Mx$. Ponadto
\[
 \|P_Mx\|^2=\sum_{i\in I}|\langle x,e_i\rangle|^2,
 \qquad
 \|x\|^2=\|x-P_Mx\|^2+
               \sum_{i\in I}|\langle x,e_i\rangle|^2,
\]
gdzie suma po dowolnym zbiorze indeksów oznacza supremum
sum skończonych. Rodzina $(e_i)_{i\in I}$ jest bazą ortonormalną
przestrzeni Hilberta wtedy i tylko wtedy, gdy dla każdego $x\in H$
zachodzi tożsamość Parsevala $\|x\|^2=\sum_{i\in I}
|\langle x,e_i\rangle|^2$; równoważnie, gdy
$M^\perp=\{0\}$.

Jeśli $(e_i)_{i\in I}$ jest bazą ortonormalną, to zbiór
$I_x=\{i:\langle x,e_i\rangle\ne0\}$ jest co najwyżej
przeliczalny. Gdy $I_x$ jest nieskończony, dla każdej bijekcji
$\nu:\N\to I_x$ zachodzi
\[
 x=\sum_{n=1}^{\infty}\langle x,e_{\nu(n)}\rangle e_{\nu(n)}
 \quad\text{w normie $H$},\qquad
 \|x\|^2=\sum_{n=1}^{\infty}|\langle x,e_{\nu(n)}\rangle|^2.
\]
Suma nie zależy od numeracji $I_x$ (zbieżność bezwarunkowa).
Jeśli $I_x$ jest skończony, obowiązuje analogiczna suma skończona.
Dla każdego skończonego $A\subseteq I$ mamy dokładny wzór na błąd
\[
 \left\|x-\sum_{i\in A}\langle x,e_i\rangle e_i\right\|^2
 =\sum_{i\in I\setminus A}|\langle x,e_i\rangle|^2.
\]
\end{twierdzenie}
\begin{proof}
Dla skończonych $A\subseteq B\subseteq I$ ortonormalność daje
$\|s_B-s_A\|^2=\sum_{i\in B\setminus A}|\langle x,e_i\rangle|^2$.
Niech $S=\sup_A\sum_{i\in A}|\langle x,e_i\rangle|^2
\le\|x\|^2$ (nierówność Bessela). Dla $\varepsilon>0$ wybieramy
$A_0$ o sumie większej niż $S-\varepsilon^2$.
Dla $A,B\supseteq A_0$ porównanie z $A\cup B$ daje
$\|s_A-s_B\|<2\varepsilon$, zatem $(s_A)$ jest siecią
Cauchy'ego. Z zupełności $H$ wynika jej zbieżność do pewnego
$s\in M$. Dla każdego $j\in I$ od chwili, gdy $j\in A$,
$\langle x-s_A,e_j\rangle=0$. Po przejściu do granicy
$x-s\in M^\perp$, więc $s=P_Mx$. Przejście do granicy
w $\|s_A\|^2=\sum_{i\in A}|\langle x,e_i\rangle|^2$
oraz twierdzenie Pitagorasa dają obie równości.
Charakterystyka bazy wynika z równoważności $M^\perp=\{0\}$
i $M=H$.

Przeliczalność $I_x$ została już niezależnie udowodniona
w Twierdzeniu~\ref{thm:countable-fourier-support13}.
Jeśli $I_x$ jest nieskończony, to zbiory
$A_N=\{\nu(1),\ldots,\nu(N)\}$ są współkońcowe w rodzinie
skończonych podzbiorów $I_x$: każdy skończony
$A\subseteq I_x$ zawiera się w pewnym $A_N$. Współczynniki
spoza $I_x$ są zerowe, więc sieć sum po podzbiorach $I$ i
sieć sum po podzbiorach $I_x$ mają tę samą granicę.
Stąd $s_{A_N}\to x$ (dla bazy $M=H$), niezależnie od $\nu$.
Suma kwadratów wynika z Parsevala. Wreszcie wektor
$x-s_A(x)$ jest prostopadły do $s_A(x)$, więc
$\|x-s_A(x)\|^2=\|x\|^2-\sum_{i\in A}|\langle x,e_i\rangle|^2$;
ponowne użycie Parsevala daje wzór na błąd.
\end{proof}

\begin{przyklad}[Różnica między rozpięciem algebraicznym i domkniętym]
\label{ex:hilbert-series-not-hamel13}
Niech $(e_n)_{n\ge1}$ będzie ciągiem ortonormalnym w $H$.
Szereg $x=\sum_{n\ge1}2^{-n}e_n$ zbiega w $H$, ponieważ
$\sum_{n\ge1}4^{-n}<\infty$. Dla każdego $n$ mamy
$\langle x,e_n\rangle=2^{-n}\ne0$, więc
$x\notin\lin\{e_n:n\ge1\}$, choć należy do domknięcia
tej podprzestrzeni. Baza ortonormalna przestrzeni Hilberta pozwala zatem na rozwinięcia
\emph{normowo zbieżne}, a baza Hamela dopuszcza wyłącznie
\emph{skończone} kombinacje liniowe.
\end{przyklad}

\begin{twierdzenie}[Istnienie baz ortonormalnych]
\label{thm:hilbert-basis13}
Każda przestrzeń Hilberta ma bazę ortonormalną (dopuszczając pustą
bazę dla $H=\{0\}$). Przestrzeń Hilberta jest ośrodkowa wtedy i tylko
wtedy, gdy posiada bazę ortonormalną co najwyżej przeliczalną.
\end{twierdzenie}
\begin{proof}
W zbiorze rodzin ortonormalnych, uporządkowanym przez rozszerzenie,
każdy łańcuch ma ograniczenie górne dane sumą rodzin.
Lemat Zorna zapewnia rodzinę maksymalną $(e_i)_{i\in I}$.
Jeśli $x\perp e_i$ dla każdego $i$ i $x\ne0$, można ją powiększyć
o $x/\|x\|$, sprzeczność. Zatem jest bazą.

Jeśli $H$ ma przeliczalną bazę, skończone kombinacje liniowe z
współczynnikami w $\mathbb Q$ (dla $\R$) lub
$\mathbb Q+i\mathbb Q$ (dla $\C$) tworzą zbiór przeliczalny i gęsty.
Odwrotnie, jeśli $H$ jest ośrodkowa, nie może zawierać
nieprzeliczalnej rodziny ortonormalnej: kule
$B(e_i,1/3)$ są rozłączne, a każdy taki otwarty zbiór zawiera
punkt ustalonego przeliczalnego zbioru gęstego.
\end{proof}


\begin{twierdzenie}[Nieskończenie wymiarowe przestrzenie Hilberta a wymiar Hamela]
\label{thm:hilbert-hamel-uncountable13}
Niech $H$ będzie nieskończenie wymiarową przestrzenią Hilberta
nad $\K\in\{\R,\C\}$. Wtedy każda baza Hamela przestrzeni $H$
jest nieprzeliczalna. Co więcej, $H$ zawiera liniowo niezależną
rodzinę mocy continuum, zatem
$\dim_{\mathrm{Hamel}}H\ge 2^{\aleph_0}$.
W szczególności twierdzenie dotyczy również ośrodkowych,
nieskończenie wymiarowych przestrzeni Hilberta, które mają
przeliczalne bazy ortonormalne.
\end{twierdzenie}
\begin{proof}
Wybierzmy nieskończony ciąg ortonormalny $(e_n)_{n\ge1}$
w $H$; istnieje on, ponieważ $H$ jest nieskończenie wymiarowa
(ewentualnie wybieramy go z bazy ortonormalnej otrzymanej
w Twierdzeniu~\ref{thm:hilbert-basis13}). Dla $t\in(0,1)$
określmy
\[
 v_t:=\sum_{n=1}^{\infty}t^{n-1}e_n.
\]
Szereg ten zbiega w normie dzięki twierdzeniu Pitagorasa i
$\sum_{n\ge1}t^{2(n-1)}=(1-t^2)^{-1}<\infty$.
Pokażemy, że rodzina $(v_t)_{t\in(0,1)}$ jest liniowo niezależna.
Niech $t_1,\ldots,t_m\in(0,1)$ będą parami różne, a
$\sum_{j=1}^m a_jv_{t_j}=0$. Biorąc iloczyn skalarny tej
równości z $e_n$ dla $n=1,\ldots,m$, dostajemy
$\sum_{j=1}^m a_jt_j^{n-1}=0$. Współczynniki $a_j$
spełniają jednorodny układ o macierzy Vandermonde'a
$(t_j^{n-1})_{1\le n,j\le m}$. Jej wyznacznik
$\prod_{1\le j<k\le m}(t_k-t_j)$ jest niezerowy, więc
$a_1=\cdots=a_m=0$. Rodzina $(v_t)_{t\in(0,1)}$ ma moc
$2^{\aleph_0}$ i jest liniowo niezależna. Każda baza Hamela
musi mieć moc co najmniej tak dużą.
\end{proof}

\begin{uwaga}[Drugi argument: twierdzenie Baire'a]
\label{rem:hilbert-hamel-baire13}
Sama nieprzeliczalność bazy Hamela wynika nawet dla każdej
nieskończenie wymiarowej przestrzeni \emph{Banacha} $X$.
Gdyby $X$ miała przeliczalną bazę Hamela $(h_n)_{n\ge1}$,
to $X=\bigcup_{n\ge1}F_n$, gdzie
$F_n=\lin\{h_1,\ldots,h_n\}$. Każda skończenie wymiarowa
podprzestrzeń $F_n$ jest domknięta (w przestrzeni unormowanej)
i ma puste wnętrze: podprzestrzeń liniowa o niepustym wnętrzu
jest całą przestrzenią, a $X$ nie jest skończenie wymiarowa.
Twierdzenie Baire'a wyklucza pokrycie zupełnej przestrzeni
metrycznej przeliczalną sumą takich zbiorów.
Argument ten dowodzi nieprzeliczalności, natomiast poprzednia
konstrukcja daje również dolne oszacowanie mocy bazy Hamela.
\end{uwaga}

\begin{uwaga}[Dwie wielkości, których nie należy mylić]
\label{rem:hamel-vs-hilbert13}
\emph{Wymiar hilbertowski} to moc bazy ortonormalnej,
a \emph{wymiar Hamela} to moc bazy algebraicznej.
Dla $H=\ell^2(\N)$ pierwszy jest przeliczalny, drugi
co najmniej continuum. Jeśli baza ortonormalna jest
nieprzeliczalna, każdy $x\in H$ nadal ma najwyżej
przeliczalnie wiele niezerowych współczynników.
\end{uwaga}

\begin{stwierdzenie}[Izometryczne współrzędne w $\ell^2(I)$]
\label{prop:hilbert-l2I13}
Niech $(e_i)_{i\in I}$ będzie bazą ortonormalną przestrzeni
Hilberta $H$. Wówczas $U:H\to\ell^2(I)$,
$Ux=(\langle x,e_i\rangle)_{i\in I}$, jest izometrycznym
izomorfizmem liniowym, gdzie
$\ell^2(I)=\{a:I\to\K:\sup_{A\subseteq I\text{ skończony}}
\sum_{i\in A}|a_i|^2<\infty\}$.
\end{stwierdzenie}
\begin{proof}
Liniowość i izometria wynikają z Parsevala. Dla dowolnego
$a\in\ell^2(I)$ sieć $\sum_{i\in A}a_ie_i$ jest Cauchy'ego:
tak jak w dowodzie Twierdzenia~\ref{thm:parseval13},
odrzucenie dostatecznie dużego skończonego zbioru współrzędnych
daje dowolnie mały ogon sumy kwadratów. Niech $x$ będzie granicą;
wtedy $\langle x,e_j\rangle=a_j$ dla każdego $j\in I$.
To dowodzi surjektywności.
\end{proof}

\begin{przyklad}[Gram--Schmidt i szeregi Fouriera]
\label{ex:gram-schmidt13}
Dla liniowo niezależnych $v_1,\ldots,v_n\in H$ konstrukcja
$u_k=v_k-\sum_{j<k}\langle v_k,e_j\rangle e_j$,
$e_k=u_k/\|u_k\|$, tworzy rodzinę ortonormalną o tym samym
rozpięciu liniowym. Dla nieskończonego ciągu liniowo niezależnego
ortogonalizujemy kolejne wyrazy i otrzymujemy bazę ortonormalną domkniętego rozpięcia tych wektorów. Dla $H=L^2([0,2\pi],\lambda)$
układ $e_n(t)=(2\pi)^{-1/2}e^{int}$, $n\in\mathbb Z$, jest
ortonormalny; jego zupełność wymaga dodatkowego twierdzenia
o gęstości wielomianów trygonometrycznych, które zostanie
omówione przy transformacji Fouriera. Nie używamy tej zupełności
bez dowodu.
\end{przyklad}

\subsection{Operatory sprzężone i reprezentacja form}

\begin{twierdzenie}[Operator sprzężony w przestrzeni Hilberta]
\label{thm:adjoint13}
Niech $H,G$ będą przestrzeniami Hilberta nad tym samym ciałem $\K$,
a $T\in\mathcal L(H,G)$. Istnieje dokładnie jeden operator
$T^*\in\mathcal L(G,H)$ taki, że
$\langle Tx,y\rangle_G=\langle x,T^*y\rangle_H$
dla wszystkich $x\in H$, $y\in G$.
Ponadto $\|T^*\|=\|T\|$, $(T^*)^*=T$ i
$(ST)^*=T^*S^*$ dla każdego $S\in\mathcal L(G,W)$,
gdzie $W$ jest kolejną przestrzenią Hilberta.
\end{twierdzenie}
\begin{proof}
Dla ustalonego $y\in G$ funkcja $x\mapsto\langle Tx,y\rangle_G$
jest ciągłym liniowym funkcjonałem na $H$. Z
Twierdzenia~\ref{thm:riesz-hilbert13} istnieje jedyny $z_y\in H$
z $\langle Tx,y\rangle_G=\langle x,z_y\rangle_H$.
Definiujemy $T^*y=z_y$. Z antyliniowości w drugim argumencie
po obu stronach, z uwzględnieniem sprzężenia skalarów, wynika
$T^*(\alpha y_1+\beta y_2)=\alpha T^*y_1+\beta T^*y_2$.
Równość norm funkcjonałów reprezentujących daje
$\|T^*y\|\le\|T\|\|y\|$, więc $T^*$ jest ograniczony.
Tożsamość dla $(T^*)^*$ wynika z hermitowskości iloczynu
skalarnego i jednoznaczności operatora sprzężonego.
Stąd $\|T\|=\|(T^*)^*\|\le\|T^*\|\le\|T\|$.
Wreszcie
$\langle STx,w\rangle=\langle Tx,S^*w\rangle
=\langle x,T^*S^*w\rangle$, co daje wzór dla iloczynu.
\end{proof}

\begin{definicja}[Formy półtoraliniowe]
\label{def:sesquilinear13}
Dla przestrzeni Hilberta $H$ nad $\K\in\{\R,\C\}$
\emph{formą półtoraliniową} nazywamy odwzorowanie
$a:H\times H\to\K$ liniowe w pierwszym i antyliniowe
w drugim argumencie. W przypadku rzeczywistym są to formy
dwuliniowe. Dla $\K=\C$ warunek antyliniowości oznacza
$a(x,\alpha y+\beta z)=\overline\alpha a(x,y)+\overline\beta a(x,z)$. Forma jest \emph{ograniczona}, jeśli
$|a(x,y)|\le M\|x\|\|y\|$ dla pewnego $M\ge0$ i wszystkich
$x,y\in H$. Najmniejsze takie $M$ oznaczamy $\|a\|$.
Formę nazywamy \emph{koercywną}, jeśli istnieje $\alpha>0$
spełniające $\operatorname{Re}a(x,x)\ge\alpha\|x\|^2$
dla wszystkich $x\in H$.
\end{definicja}

\begin{twierdzenie}[Reprezentacja ograniczonych form]
\label{thm:form-repr13}
Niech $H$ będzie przestrzenią Hilberta nad $\K$,
a $a:H\times H\to\K$ formą półtoraliniową ograniczoną.
Wówczas istnieje dokładnie jeden $A\in\mathcal L(H)$ taki, że
$a(x,y)=\langle Ax,y\rangle$ dla wszystkich $x,y\in H$.
Ponadto $\|A\|=\|a\|$.
Forma jest hermitowska ($a(x,y)=\overline{a(y,x)}$) wtedy
i tylko wtedy, gdy $A=A^*$.
\end{twierdzenie}
\begin{proof}
Dla ustalonego $x\in H$ funkcja
$y\mapsto\overline{a(x,y)}$ jest ciągłym liniowym funkcjonałem,
zatem istnieje dokładnie jeden $Ax\in H$ spełniający
$\overline{a(x,y)}=\langle y,Ax\rangle$ dla każdego $y$.
Po sprzężeniu otrzymujemy $a(x,y)=\langle Ax,y\rangle$.
Jednoznaczność daje liniowość $x\mapsto Ax$.
Równość $\|Ax\|=\sup_{\|y\|\le1}|a(x,y)|\le\|a\|\|x\|$
dowodzi ograniczoności i $\|A\|\le\|a\|$.
Z drugiej strony
$|a(x,y)|=|\langle Ax,y\rangle|\le\|A\|\|x\|\|y\|$,
stąd równość norm. Równoważność hermitowskości i $A=A^*$
wynika bezpośrednio z definicji operatora sprzężonego.
\end{proof}

\begin{twierdzenie}[Twierdzenie Laxa--Milgrama]
\label{thm:lax-milgram13}
Niech $H$ będzie przestrzenią Hilberta nad $\K\in\{\R,\C\}$,
a $a:H\times H\to\K$ formą półtoraliniową ograniczoną
i koercywną ze stałą $\alpha>0$. Dla każdego ciągłego
\emph{antyliniowego} funkcjonału $\Lambda:H\to\K$ istnieje
dokładnie jeden $u\in H$ taki, że
$a(u,v)=\Lambda(v)$ dla wszystkich $v\in H$.
Ponadto $\|u\|\le\|\Lambda\|/\alpha$, gdzie
$\|\Lambda\|=\sup_{\|v\|\le1}|\Lambda(v)|$.
\end{twierdzenie}
\begin{proof}
Niech $A\in\mathcal L(H)$ reprezentuje $a$ jak w
Twierdzeniu~\ref{thm:form-repr13}. Z koercywności i nierówności
Cauchy'ego--Schwarza mamy
$\alpha\|x\|^2\le\operatorname{Re}\langle Ax,x\rangle
\le\|Ax\|\|x\|$, zatem $\|Ax\|\ge\alpha\|x\|$.
Stąd $A$ jest iniektywny i ma domknięty obraz:
$Ax_n\to y$ implikuje, że $x_n$ jest ciągiem Cauchy'ego,
a więc $x_n\to x$ i $Ax=y$.

Jeżeli $y\perp\operatorname{ran}A$, to $A^*y=0$.
Wtedy
$\alpha\|y\|^2\le\operatorname{Re}\langle Ay,y\rangle
=\operatorname{Re}\langle y,A^*y\rangle=0$,
więc $y=0$. Zatem obraz $A$ jest gęsty i jednocześnie
domknięty, czyli $\operatorname{ran}A=H$.
Ponieważ $\Lambda$ jest antyliniowy, funkcja
$v\mapsto\overline{\Lambda(v)}$ jest ciągłym liniowym
funkcjonałem. Twierdzenie~\ref{thm:riesz-hilbert13} daje
jedyny $f\in H$ spełniający $\Lambda(v)=\langle f,v\rangle$.
Bierzemy $u=A^{-1}f$; wtedy $a(u,v)=\langle Au,v\rangle
=\Lambda(v)$. Jednoznaczność wynika z iniektywności $A$,
a oszacowanie z nierówności $\|Au\|\ge\alpha\|u\|$
i $\|f\|=\|\Lambda\|$.
\end{proof}

\begin{przyklad}[Jednowymiarowy problem brzegowy]
\label{ex:weak-poisson13}
Niech $H=\{u:[0,1]\to\C:\ u\text{ jest absolutnie ciągła},
u(0)=u(1)=0,\ u'\in L^2([0,1],\lambda)\}$,
przy czym pochodną rozumiemy prawie wszędzie.
Definiujemy $\langle u,v\rangle_H=
\int_{[0,1]}u'(t)\overline{v'(t)}\,\mathrm d\lambda(t)$.
Warunki brzegowe zapewniają dodatnią określoność. Przestrzeń $H$
jest zupełna: odwzorowanie $u\mapsto u'$ jest izometryczną bijekcją
$H\to L^2_0:=\{g\in L^2([0,1],\lambda):
\int_{[0,1]}g(t)\,\mathrm d\lambda(t)=0\}$,
którego obraz jest domknięty. Odwrotność ma wzór
$g\mapsto(x\mapsto\int_{[0,x]}g(t)\,\mathrm d\lambda(t))$.
Dla $u\in H$ nierówność Cauchy'ego--Schwarza daje
$|u(x)|\le\sqrt{x}\,\|u'\|_2$, a więc
$\|u\|_2^2\le\int_{[0,1]}t\|u'\|_2^2\,\mathrm d\lambda(t)
=\frac12\|u'\|_2^2$, a więc
$\|u\|_2\le\|u'\|_2/\sqrt2$.

Dla $f\in L^2([0,1],\lambda)$ forma
$a(u,v)=\int_{[0,1]}u'(t)\overline{v'(t)}\,\mathrm d\lambda(t)$
jest ograniczona i koercywna ze stałą $1$, zaś
$\Lambda(v)=\int_{[0,1]}f(t)\overline{v(t)}\,\mathrm d\lambda(t)$
jest antyliniowa i ciągła na $H$, ponieważ
$|\Lambda(v)|\le\|f\|_2\|v\|_2
\le\|f\|_2\|v\|_H/\sqrt2$.
Twierdzenie~\ref{thm:lax-milgram13} daje jedyne $u\in H$
spełniające $a(u,v)=\Lambda(v)$ dla wszystkich $v\in H$.
Jest to słaba postać zagadnienia $-u''=f$, $u(0)=u(1)=0$.
Aby wyznaczyć rozwiązanie jawnie, określmy
$c=\int_{[0,1]}(1-s)f(s)\,\mathrm d\lambda(s)$ oraz
\[
 g(t)=c-\int_{[0,t]}f(s)\,\mathrm d\lambda(s),\qquad
 w(x)=\int_{[0,x]}g(t)\,\mathrm d\lambda(t).
\]
Ponieważ $f\in L^2([0,1])\subseteq L^1([0,1])$,
funkcja $g$ jest absolutnie ciągła, $g'=-f$ prawie wszędzie
i $g\in L^2$. Funkcja $w$ jest absolutnie ciągła,
$w'=g$, $w(0)=0$, a z twierdzenia Fubiniego
\[
 w(1)=c-\int_{[0,1]}(1-s)f(s)\,\mathrm d\lambda(s)=0.
\]
Zatem $w\in H$ i $w''=-f$ prawie wszędzie.
Dla gładkich funkcji $v$ o zerowych wartościach na brzegu
całkowanie przez części daje
$\int_{[0,1]}w'\overline{v'}\,\mathrm d\lambda
=\int_{[0,1]}f\overline v\,\mathrm d\lambda$.
Równość rozciąga się na całe $H$ przez gęstość.
Istotnie, dla $h=v'\in L^2_0$ wybieramy
$h_n\in C^\infty([0,1])$ takie, że $h_n\to h$ w $L^2$,
i kładziemy
$\widetilde h_n=h_n-\int_{[0,1]}h_n(t)\,\mathrm d\lambda(t)$.
Wówczas $\widetilde h_n\to h$ w $L^2_0$, zaś funkcje
$v_n(x)=\int_{[0,x]}\widetilde h_n(t)\,\mathrm d\lambda(t)$
należą do $C^\infty([0,1])$, znikają na brzegu i zbiegają
do $v$ w normie $H$. Ciągłość obu stron słabej tożsamości
pozwala zatem przejść do granicy $v_n\to v$.
Jednoznaczność rozwiązania z twierdzenia Laxa--Milgrama
daje $u=w$, a więc
$u'(t)=c-\int_{[0,t]}f(s)\,\mathrm d\lambda(s)$.
\end{przyklad}

\begin{uwaga}[Droga do pochodnych Gâteaux i Frécheta]
\label{rem:gateaux-bridge13}
Jeśli $H$ jest rzeczywistą przestrzenią Hilberta, to funkcjonał
$\Phi(x)=\frac12\|x\|^2$ ma pochodną Frécheta
$D\Phi(x)h=\langle h,x\rangle$. Formalnie wynika to z
$\Phi(x+h)-\Phi(x)-\langle h,x\rangle=\frac12\|h\|^2$.
Definicję pochodnej Frécheta i jej relację do pochodnej Gâteaux
wprowadzimy w następnym rozdziale. W przestrzeni zespolonej
$\Phi$ jest funkcją rzeczywistą i rozważamy jej pochodną
\emph{rzeczywistą}, równą $h\mapsto\operatorname{Re}\langle h,x\rangle$.
\end{uwaga}
